\writeSectionFrame{\presentationTitle}{Die Schablonenmethode}
\note{
	Die Schablonenmethode bietet die Schablone eines Algorithmus
	an, die den Ablauf definiert. An bestimmten Stellen kann sich
	anderer Code einhaken und den Algorithmus beeinflussen.
}

\begin{frame}{Schablonenmethode}
	\begin{block}{Zweck}
 		Definiere das Skelett eines Algorithmus in einer Operation und delegiere einzelne Schritte an Unterklassen. Die Verwendung einer Schablonenmethode ermöglicht es Unterklassen, bestimmte Schritte eines Algorithmus zu überschreiben, ohne seine Struktur zu verändern.
 	\end{block}
\end{frame}

\begin{frame}{Allgemeines Klassendiagramm}
	\begin{center}
		\begin{tikzpicture}[scale=0.7, transform shape]
		\umlclass[type=abstract, x=5, y=0]{AbstrakteKlasse}{
		        + TemplateMethode() : void \\
		        + Methode1() : void \\
		        + Methode2() : void \\
		    }{
		        \textit{- primitiveOperation1() : void} \\
		        \textit{- primitiveOperation2() : void}
		    }
		
		    \umlclass[y=-4]{KonkreteKlasse1}{
		        + Methode1() : void \\
		        + Methode2() : void
		    }{}
		
		    \umlclass[x=10, y=-4]{KonkreteKlasse2}{
		        + Methode1() : void \\
		        + Methode2() : void
		    }{}
		
		    \umlinherit{KonkreteKlasse1}{AbstrakteKlasse}
		    \umlinherit{KonkreteKlasse2}{AbstrakteKlasse}
		\end{tikzpicture}
	\end{center}
\end{frame}

\frameWithImageOnly{\BILDERPFAD/schablone1.jpg}
\note{
	Ein Spiel läuft grundsätzlich immer gleich ab. Das Spielfeld wird aufgebaut, die Spieler machen reihum ihren
	Zug, aber die genauen Regeln des jeweiligen Spiels legen fest, wie genau das Spielfeld aufgebaut wird
	und was die Spieler in ihren Zügen machen dürfen.
}

\begin{frame}{Klassendiagramm}
	\begin{center}
			\begin{tikzpicture}[scale=0.6, transform shape]
			\umlclass[x=0, y=0]{Game}{
		        - playerCount: int \\
		        - score: int
		    }{
		        + Game(int playerCount) \\
		        + printGreeting(): String \\
		        + initializeBoard()  \\
		        + makeMove(int playerId)  \\
		        + isGameFinished(): boolean \\
		        + printWinner(): String \\
		        + playGame() \\
		        + getPlayerCount(): int
		    }
		
		    % Unterklasse PigGame
		    \umlclass[x=-4, y=-6]{PigGame}{
		        - playerScores: int[] \\
		        - currentTurnScore: int \\
		        - currentPlayer: int \\
		        - random: Random \\
		        - scanner: Scanner
		    }{
		        + PigGame(int playerCount) \\
		        + initializeBoard() \\
		        + makeMove(int playerId) \\
		        + isGameFinished(): boolean \\
		        + printWinner(): String
		    }
		
		    % Unterklasse HighLowGame
		    \umlclass[x=4, y=-6]{HighLowGame}{
		        - deck: List<Integer> \\
		        - currentCard: int \\
		        - score: int \\
		        - scanner: Scanner
		    }{
		        + HighLowGame(int playerCount) \\
		        + initializeBoard() \\
		        + makeMove(int playerId) \\
		        + isGameFinished(): boolean \\
		        + printWinner(): String
		    }
		
		    % Beziehungen (Vererbung)
		    \umlinherit{PigGame}{Game}
		    \umlinherit{HighLowGame}{Game}
		\end{tikzpicture}
	\end{center}
\end{frame}

\note{
	Unser vereinfachter Spielablauf sieht so aus, dass zuerst eine Begrüßung angezeigt wird. Das ist
	für jedes Spiel gleich. Das Initialisieren des Spielfeldes und das Ziehen ist bei jedem Spiel anders.
	Es wird auch anders herausgefunden, ob das Spiel beendet ist. Dabei wird auch der Gewinner ermittelt. Die Ausgabe
	des Gewinners funktioniert dann wieder für jedes Spiel identisch. Die Methode playGame() ist die Spielschleife,
	in der das Spiel abläuft, d.h. der Algorithmus, der von den verschiedenen Spielen angepasst werden soll. 
}

\begin{frame}[fragile]{Abstrakte Klasse}
	\begin{exampleblock}{Game.java}
 		\begin{lstlisting}
public abstract class Game {
	private int playerCount;
	
	public Game(int playerCount) {
		this.playerCount = playerCount;
	}
	
	public String printGreeting() {
		return "Willkommen zu diesem wunderbaren Spiel";
	}
	
	public abstract void initializeBoard();
	public abstract void makeMove(int playerId);
	public abstract boolean isGameFinished();
 		\end{lstlisting}
 	\end{exampleblock}
\end{frame}

\begin{frame}[fragile]{Abstrakte Klasse}
	\begin{exampleblock}{Game.java}
 		\begin{lstlisting}
public void playGame() {
    printGreeting();
    initializeBoard();
    int playersTurn = 0;
    
    while(!isGameFinished()) {
        makeMove(playersTurn);
        playersTurn = (playersTurn + 1) % playerCount;
        
    }
    
    printWinner();
}

public int getPlayerCount() {
    return playerCount;
} 		    
 		\end{lstlisting}
 	\end{exampleblock}
\end{frame}

\begin{frame}[fragile]{Konkrete Klasse 1}
	\begin{exampleblock}{PigGame.java}
 		\begin{lstlisting}
@Override
public void makeMove(int playerId) {
    // Spiellogik fuer Schweinchenspiel
}

@Override
public boolean isGameFinished() {
    for (int score : playerScores) {
        if (score >= WINNING_SCORE) {
            return true;
        }
    }
    return false;
}
 		\end{lstlisting}
 	\end{exampleblock}
\end{frame}

\begin{frame}[fragile]{Konkrete Klasse 2}
	\begin{exampleblock}{HighLowGame.java}
 		\begin{lstlisting}
@Override
public void makeMove(int playerId) {
    // Spiellogik fuer Hoch-Tief-Spiel.
    }

    if (deck.isEmpty()) {
        System.out.println("...);
        gameFinished  = true;
    }
}

@Override
public boolean isGameFinished() {
    return gameFinished;
}
 		\end{lstlisting}
 	\end{exampleblock}
\end{frame}

\begin{frame}{Umkehrung der Kontrolle}
	\begin{itemize}
 		\item Hollywood-Prinzip
 		\item Offene Rekursion
 	\end{itemize}
\end{frame}

\note{
	An dieser Stelle erkennen wir das Hollywood-Prinzip, denn die Oberklasse ruft die
	Methoden der Unterklasse auf und nicht umgekehrt, was ein wenig intuitiver wäre. 
	Das bezeichnet man auch als offene Rekursion -- Rekursion, weil dieselbe Methode
	noch einmal aufgerufen wird, und offen, weil erst zur Laufzeit entschieden wird,
	welche Implementierung tatsächlich aufgerufen wird.
}

\frameWithImageOnly{\BILDERPFAD/schablone2.jpg}
\note{
	Als nächstes verwenden wir die Schablone für einen Validierungsmechanismus. Dabei sollen die Daten von HTW- und Uni-Studenten validiert werden. Der Validierungsprozess läuft immer gleich, allerdings jeweils mit unterschiedlich aufgebauten Daten.
}

\begin{frame}[fragile]{Klassendiagramm}
\begin{tikzpicture}[scale=0.6, transform shape]

    % Klasse StudentValidation
    \umlclass[x=0, y=0]{StudentValidation}{
        - validateFirstName(student: Student): boolean \\
        - validateLastName(student: Student): boolean
    }{
        + validateStudent(student: Student): boolean \\
        \# validateMatriculationNumber(student: Student): boolean \\
        \# validateAspoYear(student: Student): boolean
    }

    % Klasse HtwStudentValidation
    \umlclass[x=-5, y=-5]{HtwStudentValidation}{
        - aspoYears: List<Integer> \\
    }{
        + validateMatriculationNumber(student: Student): boolean \\
        + validateAspoYear(student: Student): boolean
    }

    % Klasse UniStudentValidation
    \umlclass[x=5, y=-5]{UniStudentValidation}{
        - aspoYears: List<Integer> \\
    }{
        + validateMatriculationNumber(student: Student): boolean \\
        + validateAspoYear(student: Student): boolean
    }

    % Vererbungspfeile
    \umlinherit{HtwStudentValidation}{StudentValidation}
    \umlinherit{UniStudentValidation}{StudentValidation}
\end{tikzpicture}
\end{frame}

\begin{frame}[fragile]{Abstrakte Klasse}
	\begin{exampleblock}{StudentValidation.java}
 		\begin{lstlisting}
public abstract class StudentValidation {
	protected abstract boolean validateMatriculationNumber(
        Student student);
	protected abstract boolean validateAspoYear(Student student);
	
	public boolean validateStudent(Student student) {
		boolean valid = false;
		valid = validateFirstName(student);
		valid = validateLastName(student);
		valid = validateAspoYear(student);
		valid = validateMatriculationNumber(student);
		
		return valid;
	}
}
 		\end{lstlisting}
 	\end{exampleblock}
\end{frame}

\begin{frame}[fragile]{Abstrakte Klasse}
	\begin{exampleblock}{StudentValidation.java}
 		\begin{lstlisting}
private boolean validateFirstName(Student student) {
    String firstname = student.getFirstname();
    if (firstname == null || firstname.isBlank()) {
        return false;
    } else {
        return true;
    }
}

private boolean validateLastName(Student student) {
    String lastname = student.getLastname();
    if (lastname == null || lastname.isBlank()) {
        return false;
    } else {
        return true;
    }
}
 		\end{lstlisting}
 	\end{exampleblock}
\end{frame}

\begin{frame}[fragile]{Konkrete Klasse 1}
	\begin{exampleblock}{HtwStudentValidation.java}
 		\begin{lstlisting}
public class HtwStudentValidation extends StudentValidation {
	private static final int MATRICULATION_NUMBER_LENGTH = 7;
	private static final String MATRICULATION_NUMBER_START = "5";
	private List<Integer> aspoYears;
	
	public HtwStudentValidation() {
		aspoYears = new ArrayList<>();
		aspoYears.add(2016);
		aspoYears.add(2019);
		aspoYears.add(2023);
	}
 		\end{lstlisting}
 	\end{exampleblock}
\end{frame}

\begin{frame}[fragile]{Konkrete Klasse 1}
	\begin{exampleblock}{HtwStudentValidation.java}
 		\begin{lstlisting}
@Override
protected boolean validateMatriculationNumber(Student student) {
    String matriculationNumber = student.getMatriculationNumber();
    if (matriculationNumber == null || matriculationNumber.
            isBlank()) {
        return false;
    }
    
    if (matriculationNumber.
            startsWith(MATRICULATION_NUMBER_START) == false) {
        return false;
    }
    // ...
    
    return true;
}
 		\end{lstlisting}
 	\end{exampleblock}
\end{frame}

\begin{frame}[fragile]{Konkrete Klasse 2}
	\begin{exampleblock}{UniStudentValidation}
 		\begin{lstlisting}
public class UniStudentValidation extends StudentValidation {
	private static final int MATRICULATION_NUMBER_LENGTH = 10;
	private List<Integer> aspoYears;
	
	public UniStudentValidation() {
		aspoYears = new ArrayList<>();
		aspoYears.add(2017);
		aspoYears.add(2020);
		aspoYears.add(2024);
	}
 		\end{lstlisting}
 	\end{exampleblock}
\end{frame}

\begin{frame}[fragile]{Konkrete Klasse 2}
	\begin{exampleblock}{UniStudentValidation}
 		\begin{lstlisting}
@Override
protected boolean validateMatriculationNumber(Student student) {
    String matriculationNumber = student.getMatriculationNumber();
    if (matriculationNumber == null || matriculationNumber.
            isBlank()) {
        return false;
    }
    
    if (matriculationNumber.length() != 
            MATRICULATION_NUMBER_LENGTH) {
        return false;
    }
    
    return true;
}
 		\end{lstlisting}
 	\end{exampleblock}
\end{frame}

\frameWithImageOnly{\BILDERPFAD/schablone3.jpg}
\note{
	Die Schablonenmethode wird oft verwendet, um das Einhaken zu ermöglichen. Es gibt
	also Schnittstellen, um an bestimmten Stellen in einem Algorithmus etwas zu tun.
	Das wird gerne verwendet, um vor und nach einer Ausführung etwas zu tun.
}

\begin{frame}{Klassendiagramm - Beispiel}
	\begin{center}
\begin{tikzpicture}[scale=0.6, transform shape]

    % Klasse TextFileProcessor
    \umlclass[x=0, y=0]{TextFileProcessor}{
        - filePath: String
    }{
        + processFile(filePath: String): void \\
        - readFile(filePath: String): List<String> \\
        - saveFile(filePath: String, processedLines: List<String>): void \\
        \# beforeProcessing(filePath: String): void \\
        \# afterProcessing(filePath: String): void \\
        \# processLines(lines: List<String>): List<String>
    }

    % Klasse UppercaseTextFileProcessor
    \umlclass[x=-5, y=-5]{UppercaseTextFileProcessor}{
    }{
        + processLines(lines: List<String>): List<String] \\
        + beforeProcessing(filePath: String): void \\
        + afterProcessing(filePath: String): void
    }

    % Klasse LineCounterTextFileProcessor
    \umlclass[x=5, y=-5]{LineCounterTextFileProcessor}{
    }{
        + processLines(lines: List<String>): List<String]
    }

    % Vererbungspfeile
    \umlinherit{UppercaseTextFileProcessor}{TextFileProcessor}
    \umlinherit{LineCounterTextFileProcessor}{TextFileProcessor}
\end{tikzpicture}
 	\end{center}
\end{frame}

\note{
	Das nächste Beispiel ist ein Programm zur Textverarbeitung. Die Anwendung kann Textdateien einlesen, sie in einer bestimmten Weise verarbeiten (z.B. Großbuchstaben konvertieren oder Zeilen zählen), und dann das Ergebnis speichern. Dabei werden Hooks definiert, die vor und nach der Verarbeitung aufgerufen werden können. 
}

\begin{frame}[fragile]{Abstrakte Klasse}
	\begin{exampleblock}{TextFileProcessor.java}
 		\begin{lstlisting}
public abstract class TextFileProcessor {

    public final void processFile(String filePath) 
            throws IOException {
        beforeProcessing(filePath);  
        List<String> lines = readFile(filePath);
        List<String> processedLines = processLines(lines);
        saveFile(filePath, processedLines);
        afterProcessing(filePath);   
    }

    protected void beforeProcessing(String filePath) {}

    protected void afterProcessing(String filePath) {}

    protected abstract List<String> processLines(
        List<String> lines);
    // ...
}    
 		\end{lstlisting}
 	\end{exampleblock}
\end{frame}
\note {
    Der Kern des Hooking-Mechanismus sind leer implementierte Methoden, die in den Unterklassen überschrieben werden können aber nicht müssen. Sie werden zu einem bestimmten Zeitpunkt aufgerufen in dem Ablauf, den die Schablone definiert.
}

\begin{frame}[fragile]{Unterklasse ohne Hooks}
	\begin{exampleblock}{LineCounterTextFileProcessor}
 		\begin{lstlisting}
public class LineCounterTextFileProcessor 
        extends TextFileProcessor {
    @Override
    protected List<String> processLines(List<String> lines) {
        int lineCount = lines.size();
        System.out.println("...");
        return lines;
    }
}  
 		\end{lstlisting}
 	\end{exampleblock}
\end{frame}

\begin{frame}[fragile]{Hook}
	\begin{exampleblock}{UppercaseTextFileProcessor}
 		\begin{lstlisting}
public class UppercaseTextFileProcessor extends 
        TextFileProcessor {
    @Override
    protected void beforeProcessing(String filePath) {
        System.out.println("[INFO] Start ...");
    }

    @Override
    protected void afterProcessing(String filePath) {
        System.out.println("[INFO] ...");
    }

    @Override
    protected List<String> processLines(List<String> lines) {
        return lines.stream()
                    .map(String::toUpperCase)
                    .collect(Collectors.toList());
    }
 		\end{lstlisting}
 	\end{exampleblock}
\end{frame}
\note{
    Diese Klasse hakt sich in die Oberklasse ein, indem sie die beiden Methoden überschreibt.
}

\begin{frame}[fragile]{Hooks beim Sortieren}
	\begin{center}
	    \begin{tikzpicture}[scale=0.6, transform shape]
	        \umlclass[x=0, y=0]{Sorting}{
	            - numbers: List<Integer>
	        }{
	            + Sorting(numbers: List<Integer>) \\
	            + sortiere(): void \\
	            \# onBefore(): void \\
	            \# onAfter(): void \\
	            \# getNumbers(): List<Integer>
	        }
	    
	       \umlclass[x=5, y=-5]{SortOnlyEvenNumbers}{
	       }{
	           %+ SortOnlyEvenNumbers(numbers: List<Integer>) \\
	           %+ onBefore(): void
	       }
	    
	        \umlinherit{SortOnlyEvenNumbers}{Sorting}
	    \end{tikzpicture}
   \end{center}
\end{frame}

\begin{frame}[fragile]{Hooks beim Sortieren}
	\begin{exampleblock}{Sorting}
        \begin{lstlisting}
public class Sorting {
	private List<Integer> numbers;
	
	public void sortiere() {
		onBefore();
		Collections.sort(numbers);
		onAfter();
	}
	
	protected void onBefore() {	}
	
	protected void onAfter() { }
}   
        \end{lstlisting}
 	\end{exampleblock}
\end{frame}
\note{
	Der Hook-Mechanismus wird oft verwendet, um die Möglichkeit zu bieten, sich vor und nach einem
	Algorithmus einzuhaken und irgendwelche Aktionen durchzuführen. Wenn eine Unterklasse
	sich einhaken möchte, implementiert sie einfach diese beiden Methoden oder nur eine davon. Wir haben eine Klasse, die einen Algorithmus zum Sortieren einer Liste anbietet. Vor und nach dem
	Sortieren bietet sie Hook-Methoden an. 
}

\begin{frame}[fragile]{Hook}
	\begin{exampleblock}{SortOnlyEvenNumbers}
 		\begin{lstlisting}
public class SortOnlyEvenNumbers extends Sorting {
	@Override
	protected void onBefore() {
		List<Integer> zahlen = getNumbers();
		zahlen.forEach(zahl -> System.out.println(zahl));
		
		List<Integer> kopieVonZahlen = new ArrayList<>(zahlen);
		for (Integer zahl: kopieVonZahlen) {
			if (zahl % 2 != 0) {
				zahlen.remove(zahl);
			}
		}
		zahlen.forEach(zahl -> System.out.println(zahl));
	}
} 		    
 		\end{lstlisting}
 	\end{exampleblock}
\end{frame}
\note{
	Wenn eine Klasse nur gerade Zahlen sortieren möchte, überschreibt sie die onBefore()-Methode und 
	entfernt die ungerade Zahlen.
}